\begin{abstract}
% Cover navigation in complex environments is a critical challenge for autonomous robots, requiring the identification and utilization of environmental cover while maintaining efficient navigation. We propose an enhanced navigation system that enables robots to identify and utilize natural and artificial environmental features as cover, thereby minimizing exposure to potential threats. Our perception pipeline leverages LiDAR data to generate high-fidelity cover maps and potential threat maps, providing a comprehensive understanding of the surrounding environment. We train an offline reinforcement learning model using a diverse dataset collected from real-world environments, learning a robust policy that evaluates the quality of candidate actions based on their ability to maximize cover utilization, minimize exposure to threats, and reach the goal efficiently.
% Extensive real-world experiments demonstrate the superiority of our approach in terms of success rate, cover utilization, exposure minimization, and navigation efficiency compared to state-of-the-art methods. 

% Autonomous robots operating in complex environments face the critical challenge of identifying and utilizing environmental cover for covert navigation to minimize exposure to potential threats. We propose \textit{EnCoMP}, an enhanced navigation framework that integrates multi-modal perception, offline reinforcement learning, and the Adaptive Threat-Aware Visibility Estimation (ATAVE) algorithm to enable robots to navigate covertly and efficiently in diverse outdoor settings. ATAVE dynamically assesses and mitigates potential threats in real-time, enhancing the robot's ability to navigate covertly by adapting to evolving environmental and threat conditions. Our novel approach generates high-fidelity cover maps, potential threat maps, height maps, and goal maps from LiDAR point clouds, providing a comprehensive understanding of the environment. The goal map encodes the relative distance and direction to the target location, guiding the robot's navigation. We train a Conservative Q-Learning (CQL) model on a large-scale dataset collected from real-world environments, learning a robust policy that maximizes cover utilization, minimizes threat exposure, and maintains efficient navigation. Extensive experiments across diverse terrains demonstrate \textit{EnCoMP}'s superior performance compared to state-of-the-art methods, achieving a 95\% success rate, 85\% cover utilization, and reducing threat exposure to 10.5\%, while significantly outperforming baselines in navigation efficiency and robustness.



Autonomous robots operating in complex environments face the critical challenge of identifying and utilizing environmental cover for covert navigation to minimize exposure to potential threats. We propose \textit{EnCoMP}, an enhanced navigation framework that integrates offline reinforcement learning and our novel Adaptive Threat-Aware Visibility Estimation (ATAVE) algorithm to enable robots to navigate covertly and efficiently in diverse outdoor settings. ATAVE is a dynamic probabilistic threat modeling technique that we designed to continuously assess and mitigate potential threats in real-time, enhancing the robot's ability to navigate covertly by adapting to evolving environmental and threat conditions. Moreover, our approach generates high-fidelity multi-map representations, including cover maps, potential threat maps, height maps, and goal maps from LiDAR point clouds, providing a comprehensive understanding of the environment. These multi-maps offer detailed environmental insights, helping in strategic navigation decisions. The goal map encodes the relative distance and direction to the target location, guiding the robot's navigation. We train a Conservative Q-Learning (CQL) model on a large-scale dataset collected from real-world environments, learning a robust policy that maximizes cover utilization, minimizes threat exposure, and maintains efficient navigation. We demonstrate our method’s capabilities on a physical Jackal robot, showing extensive experiments across diverse terrains. These experiments demonstrate \textit{EnCoMP}'s superior performance compared to state-of-the-art methods, achieving a 95\% success rate, 85\% cover utilization, and reducing threat exposure to 10.5\%, while significantly outperforming baselines in navigation efficiency and robustness.


\end{abstract}